\section{Introduction}
\label{sec:introduction}

Dualities are one of the most pervasive organizing ideas in physics: they connect weak and strong
descriptions, high and low scales, and complementary parametrizations of the same underlying
phenomenon.
The Hamiltonian of mean force (HMF) is the exact equilibrium generator for an open quantum subsystem
\cite{gelinThermodynamicsSubensembleCanonical2009a,hiltHamiltonianMeanForce2011,talknerColloquiumStatisticalMechanics2020,rivasStrongCouplingThermodynamics2020,burkeStructureHamiltonianMean2024}.
The challenge is not definition but use: once subsystem-bath coupling is finite, physically relevant
limits are often controlled by different small parameters in different parts of the same phase plane.
In practice this makes one branch easy and its partner hard.

In this setting, the most sought-after target is typically the ground-state sector (\(\beta\to\infty\)),
which is often the least numerically forgiving.

This paper is built around a simple operational question:
\emph{when a target limit is hard in one representative, can we map it to an easier one while
keeping the same underlying state family?}
Our answer is yes, provided one works at the level of representative labels rather than
parameter-level state identities.

\paragraph{Physical framing.}
Start from
\begin{equation}
  H_{\mathrm{tot}}=H_Q+H_X+H_{\mathrm{int}},
\end{equation}
and define
\begin{equation}
  \bar{\rho}_Q(\beta)=\Tr_X e^{-\beta H_{\mathrm{tot}}}.
\end{equation}
The HMF is exactly the operator satisfying
\begin{equation}
  \frac{\bar{\rho}_Q(\beta)}{Z_X(\beta)}=e^{-\beta H_{\mathrm{MF}}(\beta)}.
\end{equation}
The key structural move is to split this into a bare thermal factor and an influence factor.
That split exposes two independent discrete freedoms:
\begin{enumerate}
  \item a strong-weak representative reflection across crossover, and
  \item an orientation-gauge choice of mean-force eigenaxis labeling.
\end{enumerate}

\paragraph{Main message.}
These maps do not claim physical-state equality under inversion.
They define a \emph{duality geometry}: they exchange asymptotic work between partner representatives.
In concrete terms, a low-\(T\), strong-influence calculation can be rewritten as a weak-influence
calculation at a dual representative, plus a separate orientation-gauge transformation for directed
observables. This maps the numerically hard ground-state limit to a numerically trivial dual regime.

\paragraph{Why this viewpoint is useful.}
The framework turns qualitative statements like ``strong coupling behaves differently'' into
explicit, testable map identities.
It also clarifies which observables are gauge-even (radial/purity sector) and which are gauge-odd
(oriented sector), which matters for fitting and collapsing data across control regimes.

\paragraph{Theory-first structure.}
The next framework sections develop representative geometry in model-independent terms.
Only then do we introduce the spin-boson qubit as Example~1, where closure formulas allow direct demonstrations:
closure-level spectral-class thresholds, dual-branch entropy formulas, and quantitative error
diagnostics.
Formal theorem-level statements and proofs are collected in the Appendix.

To avoid overclaiming, we explicitly separate closure-level representative statements from established
full-model spin-boson critical phenomenology
\cite{leggettDynamicsDissipativeTwostate1987,weissQuantumDissipativeSystems2012,bullaQuantumPhaseTransitionSubohmic2003,winterQuantumPhaseTransitionSubohmic2009}.

Relative to the closure-focused manuscript \cite{mccaulMeanForceHamiltoniansInfluence2026}, the
present work shifts emphasis from ``deriving one exact model'' to ``organizing a broad class of
mean-force calculations by dual representative structure''
\cite{jarzynskiNonequilibriumWorkTheorem2004,campisiFluctuationTheoremArbitrary2009,trushechkinOpenQuantumSystem2022}.
